\documentclass[11pt]{article}
\usepackage[margin=1in]{geometry}
\usepackage{amsmath,amssymb,amsthm}
\usepackage{booktabs}
\usepackage{natbib}
\usepackage{hyperref}
\usepackage{xcolor}

\newtheorem{theorem}{Theorem}
\newtheorem{proposition}[theorem]{Proposition}
\newtheorem{lemma}[theorem]{Lemma}
\newtheorem{corollary}[theorem]{Corollary}
\theoremstyle{definition}
\newtheorem{definition}{Definition}
\theoremstyle{remark}
\newtheorem{remark}{Remark}

\newcommand{\E}{\mathbb{E}}
\renewcommand{\P}{\mathbb{P}}
\newcommand{\R}{\mathbb{R}}
\newcommand{\pending}[1]{\textcolor{red}{[#1]}}

\title{Certifying rare overflow risk of regional pickup-and-delivery plans}
\author{Quang-Vinh Dang\\
British University Vietnam, Hung Yen, Vietnam\\
\texttt{vinh.dq4@buv.edu.vn}\quad ORCID 0000-0002-3877-8024}
\date{Draft, October 2026}

\begin{document}
\maketitle

\begin{abstract}
Medical, cold-chain and dangerous-goods distribution require that a vehicle
routing plan meets service levels of $1-\gamma$ with $\gamma$ between
$10^{-3}$ and $10^{-6}$. We study how to certify such levels for a given
plan of the vehicle routing problem with simultaneous pickup and delivery,
where each route overflows if its peak load exceeds capacity. Under gamma
demands with a Gaussian factor copula we show that, on the logarithmic scale,
a route overflows in one of two ways --- at the start, through its total
delivery, or at the end, through its total pickup --- with explicit scales
given by secular equations, and that middle stages and within-stage
combinations are close behind on long routes. This structure explains a
silent failure of single-shift importance sampling: when the weaker end
scale exceeds a quarter of the stronger one the estimator loses efficiency,
and at practical levels it underestimates by the mass of the missed mode
while reporting a small standard error. We propose a closed-form
stage-mixture estimator that places one component at the exact dominating
point of each relevant failure mode, needs no pilot simulation, and extends
to a route-day simulator with stochastic travel and service times and time
windows, where no exact method exists, and to whole hubs whose routes share
demand factors. On 148 instances built from population, postal and road data
for Northern Vietnam and Guangdong (151 million people), the estimator
matches exact dynamic programming and $10^7$-draw Monte Carlo within
statistical error, has relative variance 4--6 at every level, and costs
two to three orders of magnitude fewer samples than cross-entropy mixtures
at $10\%$ relative precision. We also report that quantum amplitude
estimation, often proposed for such tasks, does not beat this classical
estimator on these instances.
\end{abstract}

\section{Introduction}

Planning models for stochastic vehicle routing control risk through chance
constraints, robust counterparts or distributionally robust ambiguity sets
\citep{GounarisWiesemannFloudas2013,DinhFukasawaLuedtke2018,GhosalWiesemann2020}.
They evaluate route risk inside the optimiser, at moderate levels and with
exact formulas, approximations or plain sampling. Applications with strict
service levels need a different step after optimisation: a certificate that
the plan actually produced meets $\P(\text{failure}) \le \gamma$ for
$\gamma \in [10^{-6}, 10^{-3}]$, computed for every route and for the plan as
a whole, at the scale of a regional network.

At such levels crude Monte Carlo needs of the order of $100/\gamma$ draws per
route for $10\%$ relative error, which is prohibitive for thousands of routes.
Importance sampling (IS) is the standard remedy, and for Gaussian-driven
models mean shifts tuned by the cross-entropy (CE) method are the practitioner's
default \citep{deBoerEtAl2005}. The difficulty is that the failure of a
pickup-and-delivery route is a union of events --- an overflow can occur after
any stop --- and single-mode IS can then fail without warning
\citep{GlassermanWang1997}. Mixtures over failure modes
\citep{OwenMaximovChertkov2019,GeyerPapaioannouStraub2019} cure this when the
modes are known, at the price of pilot simulation when they are not.

\paragraph{Contributions.}
(i) We formulate the certification of a given plan at service levels
$10^{-3}$--$10^{-6}$, route-wise and hub-wise.
(ii) We characterise the large deviations of the peak load of a
pickup-and-delivery route under gamma marginals and a one-factor Gaussian
copula: the two end modes with explicit spectral scales, and the relative
mass of every other stage (Section~\ref{sec:structure}).
(iii) We give a sharp criterion for the failure of single-shift IS and a
diagnostic for its finite-level bias (Section~\ref{sec:single}).
(iv) We propose a pilot-free stage-mixture estimator (SMIS) built from exact
dominating points, its extension to a route-day simulator without exact
solution (MDIS), and to hubs with shared factors, with risk attribution to
routes (Section~\ref{sec:estimators}).
(v) We evaluate on a new set of 148 region-scale instances
(Section~\ref{sec:experiments}).
The importance-sampling ingredients are known --- dominating points
\citep{SadowskyBucklew1990}, multiple-path failure of single shifts
\citep{GlassermanWang1997}, mixtures as subsolutions \citep{DupuisWang2007},
factor-copula two-step schemes \citep{GlassermanLi2005,BassambooJunejaZeevi2008}
--- and we do not claim new IS theory. What is new is their exact form for the
pickup-and-delivery peak, the resulting diagnostic, and a certification
procedure that is cheap enough for regional plans.

\section{Related work}
\label{sec:related}

\emph{Risk in vehicle routing.} Chance-constrained and robust models
\citep{GounarisWiesemannFloudas2013,DinhFukasawaLuedtke2018,GhosalWiesemann2020}
handle route risk inside the optimisation with exact formulas, moment bounds or
scenario sampling. Hazardous-materials routing measures risk as accident
probability times consequence along links \citep{ErkutVerter1998}. We found
no work that certifies an already optimised routing plan at rare levels with
IS, and none that treats the mid-route peak of the simultaneous
pickup-and-delivery load \citep{Dethloff2001} as a rare event. The closest IS
precedent in logistics is stockout simulation in production--inventory
systems \citep{GlassermanLiu1996}.

\emph{Rare-event simulation.} Efficiency notions and the dominating-point
theory are surveyed by \citet{AsmussenGlynn2007}; \citet{GlassermanWang1997}
show that a single shift is inefficient when the rare set has several ways
of occurring; \citet{DupuisWang2007} build efficient schemes from
subsolutions, of which mixtures are instances; \citet{OwenMaximovChertkov2019}
bound the variance of a mixture that samples each event of a union
conditionally; CE with Gaussian mixtures handles multimodal failure regions
\citep{GeyerPapaioannouStraub2019,PapaioannouGeyerStraub2019}. For Gaussian
factor copulas, \citet{GlassermanLi2005} combine a factor shift with
conditional exponential twisting; \citet{BassambooJunejaZeevi2008} treat
extremal dependence.

\section{Problem}
\label{sec:problem}

A route visits customers $1,\dots,L$. The vehicle leaves the depot with all
deliveries $d_1,\dots,d_L$ and collects pickups $p_1,\dots,p_L$. The load after
stop $k$ is $L_k = D_k + P_k$ with $D_k = \sum_{i>k} d_i$ and
$P_k = \sum_{i\le k} p_i$ ($k=0,\dots,L$), the peak is $F = \max_k L_k$, and
the route overflows a capacity $c$ if $F > c$.

\emph{Demand model.} Each $d_i$ is gamma with shape $\kappa = 1/\mathrm{cv}^2$ and
scale $\theta^d_i = \bar d_i/\kappa$, likewise $p_i$; dependence is a
one-factor Gaussian copula per demand type: $d_i = G_{d,i}^{-1}(\Phi(Z_{d,i}))$
with $Z_{d,i} = \sqrt{\rho}\,W_d + \sqrt{1-\rho}\,E_{d,i}$, and similarly for
pickups with an independent factor $W_p$. The input
$u = (W_d, W_p, E_d, E_p)$ is standard normal. All customers of a hub share the
two factors.

\begin{definition}[Certificate]
Given $\gamma$ and a relative precision $\varepsilon$, a certificate for a
route (or a hub) is an estimate $\hat p$ of its failure probability $p$ with
$|\hat p - p| \le \varepsilon p$ at confidence $1-\delta$, together with the
check $\hat p(1+\varepsilon) \le \gamma$; for a hub, also the attribution of
$p$ to its routes.
\end{definition}

\section{Structure of the overflow event}
\label{sec:structure}

For an index set $S$ with gamma scales $\theta_S$ let
$C_S = \rho\mathbf{1}\mathbf{1}^\top + (1-\rho)I$ and
\[
\lambda(S) = \lambda_{\max}\!\left(\Theta_S^{1/2} C_S \Theta_S^{1/2}\right),
\]
the largest root of $1 = \rho\sum_{i\in S}\theta_i/(\lambda-(1-\rho)\theta_i)$.
Write $\lambda^d_k = \lambda(\{i>k\})$ for the delivery block of stage $k$,
$\lambda^p_k = \lambda(\{i\le k\})$ for its pickup block,
$\lambda_k = \max(\lambda^d_k,\lambda^p_k)$, and
$\lambda^\star = \max(\lambda^d_0,\lambda^p_L)$.

\begin{proposition}\label{prop:sandwich}
$\max(D_0, P_L) \le F \le D_0 + P_L$.
\end{proposition}

\begin{theorem}\label{thm:ld}
As $c\to\infty$ with the route fixed,
$\log\P(L_k > c) \sim -c/\lambda_k$ and $\log\P(F > c) \sim -c/\lambda^\star$.
\end{theorem}

So, on the logarithmic scale, the route overflows at the start (delivery
block) or at the end (pickup block). With
$\varrho = \min(\lambda^d_0,\lambda^p_L)/\lambda^\star$, the weaker end mode
carries relative mass $\gamma^{(1/\varrho-1)(1+o(1))}$ at level $\gamma$
(Proposition~\ref{prop:mass} in the appendix). On long routes the scales of
stages near the ends are close to $\lambda^\star$, so they also matter at
practical levels.

\begin{remark}[Pre-asymptotic regime]
At $\gamma\in[10^{-9},10^{-3}]$ and $\mathrm{cv}=0.3$ the exact dominating-point
energies are only $16$--$35\%$ of their asymptotic values $c/\lambda$ (the
log corrections of the gamma quantile are large). The scales are therefore a
qualitative guide; the estimators below use exact dominating points.
\end{remark}

\section{When a single shift fails}
\label{sec:single}

Mean-shift IS with shift $\tau$ uses the proposal $N(\tau,I)$ and weight
$w(u) = \exp(-\langle u,\tau\rangle + |\tau|^2/2)$.

\begin{theorem}\label{thm:single}
Let $\tau^\star$ be the dominating point of the dominant end mode.
(a) If $\varrho > 1/4$, the shift $\tau^\star$ is not logarithmically efficient:
$\log m_2(\tau^\star) - 2\log\P(F>c) \ge c\,(4/\lambda^\star - 1/\lambda_{\rm other})(1+o(1))$.
(b) If $\varrho \le 1/4$ and the dominant end block has its second eigenvalue
at most $\lambda^\star/4$, the shift is logarithmically efficient.
\end{theorem}

The extra second moment in (a) comes from a region the proposal almost never
visits, so the sample variance does not reveal it; in practice the estimate is
biased low by about the missed mode mass with a small reported error. Our
measurements match: on 120 real routes the single CE shift is accurate
(median error $0.1$--$0.2\%$) when $\varrho\le 0.5$ and off by $3$--$6\%$ in
median and up to $40\%$ when $\varrho > 0.75$, which is the case of balanced
deliveries and pickups (Northern Vietnam). $\varrho$ is computed in closed form
from the route data, so the diagnostic costs nothing.

\section{Estimators}
\label{sec:estimators}

\paragraph{SMIS.} For every stage $k$ and block, compute the exact minimum-norm
point $u_m$ of $\{L_k(u)\ge c\}$ (constrained minimisation started on the
block's Perron direction, using the exact quantile transform), keep the modes
whose energy $E_m=|u_m|^2/2$ is within $\log(1/\eta)$ of the smallest, and
sample from
$q = a_0 N(0,I) + \sum_m a_m N(u_m, I)$ with $a_m\propto e^{-(E_m-E_{\min})}$
and a constant defensive weight $a_0$. No pilot simulation is needed.
Theorem~\ref{thm:mixture} (appendix) gives conditions for logarithmic
efficiency.

\paragraph{MDIS.} The same construction applies to any simulator that exposes
its failure modes as scalar functions $g_m(u) > 0$. We use it for a route-day
model with lognormal travel times (shared congestion factor), service times
that grow with the handled load, and time windows; no exact method applies.

\paragraph{Hubs.} The routes of a hub share the factors, so the hub failure
$\bigcup_r\{F_r > c\}$ is again a union of stage events. We embed the route
modes in the hub input, compute the likelihood ratio with sparse inner
products (each component touches the two factors and one route's
coordinates), and estimate from the same draws the probability that each
route fails and that it is the only one failing --- an attribution of the hub
risk to routes. Hubs of a region have independent factors, so the region risk
is $1-\prod_h(1-p_h)$.

\section{Experiments}
\label{sec:experiments}

\subsection{Instances}
We built 148 instances, one per depot area: the 25 provinces of Northern
Vietnam (Hanoi split into two hubs) and 122 county-level areas of Guangdong.
Service points aggregate Kontur population hexagons (15\,000 people per point
in Vietnam, 4\,000 in Guangdong) by weighted $k$-means; distances are OSRM road
distances. Daily demand rates come from official postal statistics: 0.087
parcels per person per day in Vietnam (3.2 billion items in 2024) and, for
Guangdong, 0.34 received and 0.91 sent (national and provincial 2024
volumes); 1~kg per parcel and truck capacities of 15 and 30 tonnes are
assumptions. The set has 40\,099 service points serving 151 million people.
Plans are solved with the deterministic and the sample-average chance gates;
we certify the latter.

\subsection{Routes under the demand model}
On 120 routes (stratified by region and length) and
$\gamma\in\{10^{-3},10^{-4},10^{-5},10^{-6}\}$, the CE mixture used as a
classical rival agrees with exact dynamic programming (fine quadrature over
both factors, routes with $L\le 8$) in 178 of 178 cases, maximal deviation
$0.74\%$. SMIS has relative error at most $0.6\%$ (90\% quantile) in every
$\varrho$ band, relative variance $3.8$--$5.9$ (CE: $3.8$--$6.0$), $2$--$7$
components, and is built in $0.4$--$0.9$\,s. Including CE's pilot draws
($0.8$--$1.2\times10^5$ per route), CE costs $224$--$276$ times more samples
than SMIS at $\varepsilon=0.1$ and $3.2$--$3.8$ times more at $\varepsilon=0.01$.

\subsection{Route days without an exact method}
On 60 routes with thresholds that split the failure probability evenly
between load and time modes, MDIS deviates from $10^7$-draw crude Monte Carlo
by $0.76\%$ in median ($1.77\%$ at the 90\% quantile, within the Monte Carlo
error of about $1\%$) at $\gamma=10^{-3}$; CE deviates by $0.89\%$ and
$1.75\%$; the single CE shift by $2.2\%$ and $22\%$. MDIS and CE agree within
$0.20$--$0.24\%$ at all levels, while the single shift diverges by up to
$28\%$, $63\%$ and $98\%$ (90\% quantile) at $10^{-3}$, $10^{-4}$ and
$10^{-5}$. Relative variances: MDIS $4.3$--$6.1$, CE $4.4$--$6.0$, single shift
$357$--$1\,430$. MDIS is built in $3$--$4$\,s against $43$--$55$\,s for CE, and
CE costs $1\,750$--$1\,830$ times more samples at $\varepsilon=0.1$ and
$18$--$19$ times more at $\varepsilon=0.01$.

\subsection{Hubs}
On 66 hubs (all of Northern Vietnam, 40 sampled from Guangdong), the hub
estimate lies between $\max_r p_r$ and $\sum_r p_r$ in every case and agrees
with $2\times10^6$-draw crude Monte Carlo on the 19 smallest hubs
($|z|\le 1.67$). Ignoring the shared factors --- treating routes as
independent --- overstates hub risk by a factor $1.75$--$2.2$ in median and up
to $3.2$, and $27$--$37\%$ of failing days have two or more failing routes:
standby capacity should be planned for clusters of failures.
\pending{Capacity buffers for $\gamma=10^{-3}$ and $10^{-5}$ and region totals
from the calibrated rerun.}

\subsection{A note on quantum amplitude estimation}
Quantum amplitude estimation offers a quadratic reduction in the number of
simulator calls for estimating a probability \citep{Brassard2002,Montanaro2015}.
On the same 120 routes, with explicit constants for a rigorous schedule, it
needs $0.003$--$0.5$ times the calls of crude Monte Carlo, but $2.5$--$380$
times more calls than classical IS once IS is available, and with an IS-equipped
oracle it is cheaper than classical IS only in an idealised accounting at
$\varepsilon\le 3\%$. On this problem the structure exploited by the mixture,
not the quadratic speed-up, is what makes certification cheap.

\section{Conclusion}
\label{sec:conclusion}

Rare overflow of a pickup-and-delivery route has a simple anatomy --- two end
modes with closed-form scales, near-end stages close behind --- and this
anatomy decides when common IS shortcuts fail and how to build an estimator
that needs no tuning. The resulting certificates are accurate to a fraction
of a percent at $10^{-3}$--$10^{-6}$, extend to route days without exact
solution and to hubs with shared demand factors, and are cheap enough to run
on regional networks. Open points: full proofs of the efficiency results at
the stated generality, calibration of the capacity search, region-wide
factors that couple hubs, and data-driven demand and travel models.

\appendix
\section{Proofs}
\label{app:proofs}

\paragraph{Proof of Proposition~\ref{prop:sandwich}.} $L_0 = D_0$ and $L_L = P_L$
are stage loads, and $D_k\le D_0$, $P_k\le P_L$ for all $k$. \qed

\paragraph{Lemma A1 (quantile).} For the gamma law,
$x(z) = G^{-1}(\Phi(z)) = \theta z^2/2 + \theta(2\kappa-1)\log z + O(1)$ as
$z\to\infty$, and $0\le x(z)\le x(0)$ for $z\le 0$.

\paragraph{Proof sketch of Theorem~\ref{thm:ld}.} Upper bound: by Lemma A1,
$D_k\le \tfrac12 Z^\top\Theta Z + R$ with $R$ a sum of logarithmic terms; for
$t<1/\lambda$, $\E\exp(t\,\tfrac12 Z^\top\Theta Z) =
\det(I - t\Theta^{1/2}C\Theta^{1/2})^{-1/2}<\infty$, so Chernoff's bound gives the
rate $1/\lambda$, and $R$ is negligible. The delivery and pickup blocks have
disjoint inputs, so their moment generating functions multiply and $L_k$ has
rate $1/\lambda_k$. Lower bound: a ball of radius $O(1)$ around
$(1+\eta)\sqrt{2c/\lambda}\,u^\star$, with $u^\star = B^\top\Theta^{1/2}v/\sqrt\lambda$
($Z=Bu$, $BB^\top = C$, $v$ the Perron vector), lies in the event. For $F$ use
Proposition~\ref{prop:sandwich}. \qed

\begin{proposition}\label{prop:mass}
For a fixed gap $\lambda_k<\lambda^\star$,
$\P(L_k>c)/\P(F>c) = \exp(-c(1/\lambda_k-1/\lambda^\star)(1+o(1)))$.
\end{proposition}

\paragraph{Lemma A2.} Let $h(u) = \tfrac12 u^\top M u$ with $M=A^\top\Theta A$,
$A\ge 0$, eigenvalues $\lambda_1>\lambda_2$, $\kappa=\lambda_1/\lambda_2\ge 4$,
$h_+(u) = \tfrac12|\Theta^{1/2}(Au)_+|^2$, and $\tau = r v_1$ with
$r^2 = 2c/\lambda_1$. Then $h_+(u)\ge c$ implies $|u+\tau|^2\ge 4|\tau|^2$.
\emph{Proof.} Write $u = a v_1 + u_\perp$, $b=|u_\perp|$. If $a\ge0$, then
$c\le\tfrac12(\lambda_1a^2+\lambda_2b^2)$, so $b^2\ge\kappa(r^2-a^2)$ and
$|u+\tau|^2\ge(1-\kappa)a^2+2ar+(1+\kappa)r^2$, a concave function of $a$
whose minimum on $[0,r]$ is $\min((1+\kappa)r^2,4r^2) = 4r^2$ for $\kappa\ge3$;
for $a\ge r$, $|u+\tau|^2\ge(a+r)^2\ge4r^2$. If $a<0$, then $Av_1\ge0$ gives
$(Au)_+\le(Au_\perp)_+$, so $b^2\ge\kappa r^2\ge4r^2$. \qed

\paragraph{Proof of Theorem~\ref{thm:single}(a).} $\tau^\star$ acts on the
dominant block's inputs only and the other end event depends on the other
block's inputs only, so $m_2\ge\E[w]\,\P(\text{other end})=
e^{|\tau^\star|^2}\P(\text{other end}) = \exp(2c/\lambda^\star - c/\lambda_{\rm other}+o(c))$,
while $2\log\P(F>c) = -2c/\lambda^\star(1+o(1))$. \qed
(b) follows from Proposition~\ref{prop:sandwich} and Lemma A2 applied to the
block-diagonal form of the two end blocks; the logarithmic remainders are
handled by Cauchy--Schwarz (sketch).

\begin{theorem}[sketch]\label{thm:mixture}
With a constant defensive weight, SMIS is logarithmically efficient if every
retained block has $\kappa\ge4$ and in every retained stage one block
dominates the other by a factor of at least 4; stages with
$\lambda_k\le\lambda^\star/2$ need not be retained.
\end{theorem}

\bibliographystyle{plainnat}
\bibliography{references}
\end{document}
